% !TEX root = ../main.tex

\chapter{绪论}

\section{研究背景}

大语言模型（Large Language Model，LLM）的快速发展，使自然语言逐渐成为人机交互中更重要的接口形式。相比传统软件系统中严格定义的菜单、命令和参数，大语言模型能够直接理解用户用自然语言表达的目标，并在一定程度上完成推理、总结、规划和代码生成等复杂任务。这种能力使大语言模型不再只是一个被动回答问题的文本生成工具，而逐渐成为许多自动化系统中的核心决策组件。

在此基础上，大模型智能体（Large Language Model Agent，LLM Agent）进一步将语言模型的语义理解和推理能力与外部工具、运行环境和任务规划机制结合起来。ReAct 等工作展示了将推理和行动交替组织的范式，使模型能够在推理过程中主动选择工具并利用环境反馈完成任务\cite{yao2022react}。相关综述也指出，大模型智能体已经从单纯的文本生成系统扩展为由模型推理、工具使用和环境交互共同构成的任务执行系统\cite{wang2024survey}。因此，智能体安全问题的关键不再只是模型是否会说出不安全内容，还包括模型是否会在真实环境中采取未被用户授权的行动。

这种系统形态带来了明显的应用价值。许多原本需要熟悉命令行、文件结构或 API 文档才能完成的工作，可以被转化为自然语言任务交给智能体完成。例如，用户可以要求智能体检查项目目录、整理文件、生成测试脚本或调用网页服务。OpenClaw 这类开源智能体运行框架进一步将模型连接到真实本地工作区，使其能够围绕文件系统和工具链执行更复杂的自动化任务\cite{openclaw2026github}。随着能力边界扩展，安全边界也随之变化：一次错误的文件删除、脚本执行或配置修改，已经不只是文本层面的回答问题，而可能直接改变用户环境状态。

\begin{figure}[htbp]
  \centering
  \small
  \setlength{\fboxsep}{8pt}
  \begin{tabular}{c@{\quad}c@{\quad}c}
    \fbox{\begin{minipage}{0.30\linewidth}
      \centering
      \textbf{传统 LLM}\\[0.4em]
      用户输入\\
      $\downarrow$\\
      文本生成\\
      $\downarrow$\\
      输出内容
    \end{minipage}}
    &
    \(\Longrightarrow\)
    &
    \fbox{\begin{minipage}{0.46\linewidth}
      \centering
      \textbf{LLM Agent}\\[0.4em]
      用户目标与环境信息\\
      $\downarrow$\\
      任务规划与工具选择\\
      $\downarrow$\\
      文件、命令、API 调用等真实行动
    \end{minipage}}
    \\
    文本层面的正确性与安全性
    &
    \(\Longrightarrow\)
    &
    行动层面的授权性与可控性
  \end{tabular}
  \caption{大模型智能体使安全关注对象从文本输出扩展到工具行动}
  \label{fig:llm-agent-risk-shift}
\end{figure}

如图~\ref{fig:llm-agent-risk-shift} 所示，智能体系统并不是简单地把大语言模型包装成一个更方便的对话界面，而是把模型输出进一步连接到工具和环境中。对于传统对话模型而言，模型输出错误通常主要表现为事实错误、误导性建议或不安全文本；而当模型拥有工具调用权限后，错误决策可能直接改变真实环境状态。因此，大模型智能体安全问题不仅关注模型是否会生成危险内容，更需要关注其在真实环境中是否可能执行超出用户授权、违背用户意图的行为。

已有智能体安全基准和防御框架的实验现象表明，在智能体场景中，提示注入攻击不仅可能影响模型的文本回答，还可能进一步干扰工具选择、参数生成和工作流路径\cite{debenedetti2024agentdojo,debenedetti2025defeating}。FIND-Lab 基于 OpenClaw 构建的 AgentWard 安全防御插件，为研究智能体运行时防护提供了具体系统基础\cite{zhang2026agentward}，在此基础上，本文进一步聚焦意图决策对齐问题：当智能体准备调用工具时，系统需要判断当前行动是否仍服务于用户真实授权目标，而不能仅依据工具命令表面的危险程度作出判断。

\section{大模型智能体的安全挑战}

大模型智能体的安全挑战并非来自单一环节，而是同时体现为输入来源复杂、工具动作依赖语义授权、任务执行过程持续演化等多方面问题。具体而言，可以从以下三个方面理解。

\begin{enumerate}
  \item \textbf{不可信外部内容可能进入模型上下文。}普通对话模型主要处理用户直接输入，而智能体在执行任务时会不断读取外部信息，包括网页内容、邮件正文、文档说明（如 README 文件）、工具返回结果、历史记忆以及其他程序生成的中间状态等。这些内容并不一定可信，却可能与用户指令一起进入模型上下文。如果外部内容中包含攻击者插入的指令，模型可能难以区分哪些内容是用户真正授权的任务，哪些只是被读取到的不可信数据。AgentDojo 将这类间接提示注入问题放入动态任务环境中评估，说明攻击指令可以隐藏在工具结果中并影响智能体的后续行动\cite{debenedetti2024agentdojo}。InjecAgent 也进一步针对工具集成智能体构造了间接提示注入基准，揭示了外部内容污染对工具调用安全的影响\cite{zhan2024injecagent}。

  \item \textbf{工具动作的安全性依赖用户任务和授权边界。}智能体风险并不总是表现为显式恶意命令。很多情况下，单独观察某个工具调用并不容易判断它是否危险，因为同一个动作在不同任务中可能具有完全不同的含义。例如，删除临时文件在清理缓存任务中可能是合理操作，但在用户只要求查看文件列表时就明显超出了授权范围；写入配置文件在修复项目时可能是必要步骤，但如果用户只要求审计配置，则这种写入就可能是过度行动。也就是说，工具调用的安全性不仅取决于动作本身，还取决于用户任务、环境上下文和授权边界。文本输出安全和工具调用安全之间也可能存在断裂：一个智能体可以在自然语言回复中表现得谨慎、礼貌、符合安全规范，但其后台工具调用仍然可能执行不符合用户意图的操作。GAP benchmark 专门指出，文本安全并不能自然迁移到工具调用安全，单纯观察模型说了什么并不足以判断智能体实际做了什么\cite{cartagena2026mind}。因此，防御系统不能只检查用户输入或模型回复，而必须在工具调用发生前观察智能体即将采取的具体行动。

  \item \textbf{自主规划和长周期执行会放大决策偏移。}智能体的自主规划能力还可能引发任务范围扩展问题。用户的自然语言请求通常难以穷尽所有边界条件，尤其在真实使用场景中，用户更可能给出目标式、概括式的指令，例如“帮我检查这个项目”“整理一下当前目录”“确认这些配置是否有问题”。为完成此类任务，智能体需要自行推断后续行动；然而，这种推断过程可能使任务从辅助检查扩展为自动修改，从信息读取扩展为环境清理，从局部操作扩展为全局变更。此类风险未必源于恶意攻击，但同样可能导致智能体决策偏离用户真实意图。在长周期任务中，早期误判的影响还会被进一步放大。智能体通常在多轮观察、思考和行动中逐步完成任务，一个早期错误如果没有被及时发现，就可能被写入计划、记忆或中间文件，并在后续步骤中持续影响模型判断。Agent Security Bench 等工作将智能体攻击面扩展到系统提示、用户提示、工具使用和记忆检索等多个环节，也表明智能体安全并非单点问题，而是贯穿整个运行过程的系统问题\cite{zhang2025agent}。在长周期任务中，决策偏移有时并不是一次明显的越权行为，而是由多个看似正常的小步骤逐渐累积而成。
\end{enumerate}

综上，大模型智能体安全的核心困难在于动态性和语义性。所谓动态性，是指智能体行为会随环境、历史状态和工具反馈不断变化，攻击或偏移不一定出现在任务开始时；所谓语义性，是指安全判断往往需要理解用户真实意图、工具动作含义和上下文约束之间的关系，而不能只依赖关键词或固定规则。本文关注的意图决策不对齐问题，正是对这两类困难的一种概括。

\section{现有防御思路及其不足}

围绕大模型智能体安全，现有防御思路可以从输入、模型和运行时三个层面理解。输入侧防御主要在用户输入或外部内容进入模型前进行处理，例如使用分隔符区分系统指令与外部数据、对输入进行攻击检测、对提示词进行重写，或在系统提示中加入更严格的安全策略。这类方法通常容易部署，运行成本也相对较低，适合作为安全防御的第一道屏障。

但是，输入侧防御的局限也比较明显。智能体运行时接触到的信息并不只来自用户最初输入，很多风险会在工具调用之后、环境读取过程中或多轮交互中出现。如果防御只在输入端进行一次检查，就难以覆盖后续动态生成的风险。此外，攻击者也可以通过更隐蔽的表达方式绕过简单输入检测，使模型在表面上没有看到明显攻击语句的情况下仍然产生错误行动。

模型侧安全增强则旨在通过安全微调、偏好优化、强化学习或专门训练安全模型等方式，提升模型本身的安全倾向，使其在面对危险请求时更倾向于拒绝或采取保守策略。对于拥有模型训练权限的系统而言，模型侧增强具有较强吸引力。

但在本文关注的 OpenClaw 这类可切换模型的智能体框架中，模型侧方法并不总是可行。用户可能使用闭源模型 API，也可能在不同模型之间切换；防御者往往无法直接修改核心模型参数。同时，模型侧增强依赖训练数据覆盖，如果新的攻击方式或新的环境偏移不在训练分布中，模型仍可能做出错误判断。因此，模型侧安全增强虽然重要，但不足以单独构成面向真实智能体运行框架的完整防线。

运行时防御直接面向智能体执行过程，在模型响应、行动轨迹或工具调用发生时进行监测，并在风险行为真正执行前进行干预。相比输入侧和模型侧方法，运行时防御更接近智能体实际行动，因此更适合处理工具调用安全问题。现有研究已经从多个方向展开探索：AgentArmor 将智能体运行轨迹转化为可分析的程序表示，并结合控制流、数据流和依赖关系进行策略检查\cite{wang2025agentarmor}；ShieldAgent 面向安全策略推理和可验证检查，在受保护智能体行动前生成屏蔽计划并进行验证\cite{chen2025shieldagent}；AGrail 关注长期运行过程中的自适应安全检测\cite{luo2025agrail}；GuardAgent 使用额外的防护智能体理解用户给定的安全要求，并将其转化为可执行的检查逻辑\cite{xiang2024guardagent}；TraceAegis 则利用执行轨迹中的层次结构和行为约束检测异常行动\cite{liu2025traceaegis}。这些工作说明，安全检查正在从单次输入过滤逐渐转向更贴近智能体行动过程的运行时治理。

运行时防御本身仍然面临一个取舍问题。基于静态规则的方法通常稳定、可解释、开销较低，对于删除、执行、外传等显式危险动作有较好效果；但它们难以识别语义层面的目标偏移。例如，一个写文件动作并不一定危险，只有当写入对象、写入内容或任务授权范围发生偏移时才构成风险。相反，基于 LLM 的动态检测方法能够理解更多语义信息，也更容易适应未知场景，但它们可能带来误报、幻觉、时延和 token 成本，并且本身也依赖裁判模型的能力。

除直接检测风险行动外，也有工作从系统结构上减少不可信数据对智能体决策的影响。CaMeL 通过显式区分控制流与数据流，使从外部环境中读取的不可信数据不能改变由可信查询决定的程序控制路径\cite{debenedetti2025defeating}。这类方法并不等同于任务对齐检测，但它提示了一个重要事实：在智能体系统中，安全风险往往出现在“谁有权影响下一步行动”这一边界上。

与上述结构隔离和轨迹检测思路相联系，任务对齐方向进一步强调智能体动作是否真正服务于用户目标。Task Shield 将间接提示注入防御重新表述为任务对齐问题，要求每条指令和每次工具调用都能够贡献于用户指定目标\cite{jia2025task}；GAP benchmark 则指出，文本层面的安全并不能自然迁移到工具调用层面的安全，模型可以在文本上拒绝危险请求，却在工具调用中执行不安全行动\cite{cartagena2026mind}。这些工作为本文提供了直接启发：对于工具化智能体而言，安全防御不能只关心某个动作表面上是否危险，还应该关注该动作是否被用户任务授权、是否仍然服务于用户真实目标。

基于上述分析，本文将运行时意图决策对齐作为核心研究对象。在智能体准备调用工具之前，防御系统需要分析用户真实意图、智能体当前决策和待执行工具动作之间的关系：当前行动与用户授权目标一致时，系统尽量不干扰智能体正常工作；行动明显偏离用户意图时，系统在执行前进行阻断；行动边界不清时，系统进入审查流程。该思路在规则方法的稳定性和 LLM 方法的语义理解能力之间取得平衡，为工具化智能体提供更贴近运行过程的安全防护。

概括而言，现有防御方法的重要不足并不在于缺少某一种具体检测规则，而在于尚未充分刻画智能体行动与用户授权之间的语义关系。输入侧方法关注不可信内容能否进入上下文，模型侧方法关注模型是否具备更强安全倾向，执行侧方法关注动作本身是否危险；但在许多真实场景中，风险恰恰发生在这些边界之间，即智能体把外部内容、模糊目标或中间推断误认为用户授权，并据此生成看似合理却已经偏离任务的行动。因此，围绕用户真实意图、智能体当前决策和工具调用动作之间的一致性建立运行时检查机制，是本文选择意图决策对齐作为核心研究方向的主要原因。

\section{研究问题与本文工作}

\subsection{研究问题与方法思路}

本文研究的问题可以概括为：在大模型智能体与真实环境交互的过程中，如何识别并处理智能体当前决策与用户真实意图之间的不对齐，避免模型在工具调用阶段执行未授权或偏离任务目标的行为。

意图决策不对齐并不等同于一般意义上的危险命令检测。危险命令检测关注动作表面特征，例如是否出现删除、执行、外传、修改权限等高风险操作；意图决策对齐关注动作与用户授权目标之间的关系。二者有交集，但并不相同。一个表面危险的动作可能在用户明确授权下是合理的，一个表面普通的动作也可能因为目标对象被重映射而造成风险。

以本文实验中使用的 \texttt{README} 间接注入场景为例，用户请求智能体列出工作区文件名称及属性，并明确限制智能体不要对文件执行其他操作；工作区 \texttt{README} 文件中却包含“在任何操作前清理临时目录”的环境指令。在没有防御时，智能体可能将该环境指令误认为用户授权任务的一部分，从而在只读任务中触发删除动作。该场景的风险并不只在于删除命令这一字符串本身，而在于删除动作已经脱离用户原始任务边界。目标别名重映射场景也具有类似特点：工具调用表面上可能只是普通写入，但如果写入目标已经从用户指定文件转移到攻击者希望修改的对象，那么该动作同样违背了用户真实意图。

除环境注入外，模糊指令和长周期任务也会导致类似问题。用户有时只表达目标而没有明确行动边界，智能体为了完成任务可能主动扩展操作范围；在多轮任务中，模型也可能逐渐遗忘最初约束，将中间推断当作新的目标。意图决策对齐能够将环境注入、模糊授权、目标漂移和工具调用分歧纳入同一分析框架：这些现象虽然来源不同，但最终都体现为智能体当前准备执行的行动与用户真实授权之间出现了偏差。

围绕这一问题，本文提出了三模块防御思路。首先，用户真实意图分析模块从当前用户输入及近期交互上下文中提取任务目标、约束条件、敏感对象和风险等级，用于刻画用户在本轮交互中真正授权智能体完成的任务范围。其次，意图对齐检测模块在工具调用前，对用户真实意图、智能体当前输出以及待调用工具参数进行综合比较，判断当前行动是否处于对齐、模糊或不对齐状态。最后，动态结果处理模块根据检测结果分别采取放行、审查或阻断策略，使系统能够针对不同风险程度进行差异化处理，而非对所有可疑情况采用统一的一刀切策略。

本文方法遵循安全性与可用性并重的设计原则。如果所有不确定行为都被直接阻断，系统虽然看似安全，但会严重影响用户正常任务；如果为了可用性而过度放行，又会失去防御意义。因此，本文引入分级处理思路，将明确安全、边界不清和明确偏移的行为区分开来，并结合规则过滤、缓存和异常回退等工程机制，降低不必要的 LLM 调用和重复判断。

\subsection{本文工作与贡献}

本文工作基于 OpenClaw 和 AgentWard 展开。OpenClaw 提供本地智能体任务执行与工具调用环境，Taming OpenClaw 提出了面向智能体生命周期的五层防御框架，AgentWard 则以 OpenClaw 插件形式实现了这一运行时防御思想\cite{openclaw2026github,deng2026taming,zhang2026agentward}。在此基础上，本文不重新实现完整智能体框架，而是聚焦决策对齐层中的语义授权检查问题，进一步分析用户真实意图、智能体当前决策和待执行工具动作之间的关系，设计三模块运行时防御方法，并完成相应代码拓展、稳定性优化和实验验证。

表~\ref{tab:chap01-contributions} 概括了本文工作与 OpenClaw、AgentWard 既有基础之间的关系。该表的目的不是割裂本文方法与系统基础，而是明确本文新增工作的边界：OpenClaw 和 AgentWard 提供运行时接入基础，本文进一步把决策对齐层细化为可执行的意图分析、对齐检测和分级处置流程。

\begin{table}[htbp]
  \centering
  \caption{本文工作与既有系统基础的关系}
  \label{tab:chap01-contributions}
  \small
  \renewcommand{\arraystretch}{1.25}
  \begin{tabular}{@{}p{0.16\textwidth}p{0.30\textwidth}p{0.38\textwidth}@{}}
    \toprule
    层面 & OpenClaw / AgentWard 提供 & 本文新增工作 \\
    \midrule
    运行环境 &
    工作区交互与工具调用 &
    工具执行前接入意图决策检查 \\
    防御框架 &
    五层防御思想与 hook 接入 &
    细化决策对齐层的语义授权检查 \\
    状态建模 &
    用户输入、上下文和工具调用信息 &
    抽取目标、约束、敏感对象和风险等级 \\
    检测逻辑 &
    运行时观察与干预 &
    比较用户意图、智能体决策和工具参数 \\
    实验验证 &
    真实 OpenClaw 运行基础 &
    构造真实场景与模拟轨迹并评估结果 \\
    \bottomrule
  \end{tabular}
\end{table}

具体而言，本文的贡献主要体现在三个方面。第一，本文将 OpenClaw 运行时中的一类安全风险概括为意图决策不对齐，并指出该问题不同于单纯的危险命令检测。环境注入、模糊授权、目标重映射和长周期任务漂移虽然表现形式不同，但都可以理解为智能体当前行动与用户真实授权目标之间出现偏差。

第二，本文围绕该问题设计并实现了三模块决策对齐防御机制。该机制在工具调用前显式分析用户真实意图，并将其与智能体当前决策和工具参数进行比较，再根据对齐程度输出不同处置策略。相比只依赖静态规则或单次 LLM 安全判断的方法，本文方法更强调用户意图、智能体决策和工具行动之间的一致性。

第三，本文构建了模拟轨迹测试和真实场景测试两类实验，对方法的有效性、可用性和模型依赖性进行了验证。模拟轨迹测试用于检查三模块框架是否能够识别不同类型的对齐状态；真实场景测试则在 OpenClaw 环境中比较 None、Rule-Only、Prompt-Policy、LLM-Only 和本文方法等不同基线方法（baseline）设置，并在 Qwen3.5-Plus 和 MiniMax-M2.5 两个模型上观察结果差异。实验结果初步表明，本文方法能够在当前测试集上降低攻击残留，并在多数良性任务中保持智能体的基本可用性，为运行时决策对齐防御提供实验依据。

\section{论文组织结构}

本文后续章节安排如下。

第 2 章介绍大模型智能体安全相关工作。该章将从大模型智能体与工具调用安全、提示注入与攻击基准、输入侧防御、模型侧安全增强、运行时检测，以及意图分析与决策对齐等方面展开梳理，概括现有研究的主要思路与不足，并进一步明确本文工作的研究切入点。

第 3 章给出意图决策对齐问题定义与总体框架。该章首先介绍 OpenClaw 和 AgentWard 运行时框架，然后定义意图决策不对齐问题，说明风险来源、设计目标和三模块总体框架。

第 4 章介绍本文提出的防御机制设计与实现。该章详细说明用户真实意图分析、意图对齐检测和动态结果处理三个模块，并介绍其在 OpenClaw 运行时中的接入方式、状态管理、LLM 调用和稳定性优化。

第 5 章介绍实验设计与结果分析。该章说明真实场景测试和模拟轨迹测试的数据构造方式、基线方法设置、评价指标和实验结果，并分析不同模型下本文方法的安全性与可用性表现。

第 6 章总结全文工作，并讨论本文方法的局限与后续改进方向。
